\documentclass[10pt,a4paper]{amsart}
\usepackage[margin=19mm]{geometry}
\usepackage{amsmath,amssymb,mathtools}
\usepackage[scr=rsfs,cal=euler]{mathalfa}
\usepackage{xfrac}
\usepackage[unicode,hidelinks,bookmarks=false]{hyperref}
\newcommand{\ie}{i.e.\ }
\newcommand{\cf}{cf.\ }
\newcommand{\resp}{respectively\ }
\newcommand{\Subsection}{\subsection}
\providecommand{\nobreakdash}{\nobreak\hbox{-}}
\setlength{\emergencystretch}{2em}
\multlinegap0pt
\usepackage{tikz}
\usepackage{amssymb}
\usepackage[shortlabels]{enumitem}
\usepackage{algorithm}\usepackage{algpseudocode}
\usepackage{stackengine}
\usepackage{dsfont}
\usepackage{xcolor}
\usepackage[normalem]{ulem}
\newtheorem{lemma}{Lemma}
\newtheorem{theorem}{Theorem}
\newcommand{\sq}{\mathbin{\square}}
\title{A cyclic flat embedding theorem for transversal $q$-matroids}
\author{Andrew Fulcher}
\date{Source: arXiv:2603.13550v1 (CC BY 4.0)}
\begin{document}
\maketitle

\noindent\emph{Source and licence.} A cyclic flat embedding theorem for transversal $q$-matroids. Version arXiv:2603.13550v1 (CC BY 4.0), distributed by arXiv under the Creative Commons Attribution 4.0 International licence, http://creativecommons.org/licenses/by/4.0/, as recorded in the arXiv licence metadata for this exact version and shown on the abstract page https://arxiv.org/abs/2603.13550v1. Full licence notice: \texttt{licenses/arxiv-2603.13550-CC-BY-4.0.txt}.\par\smallskip

\noindent\textit{Editorial scope.} Counted unit: the article's theorem thm:free\_product\_transversal (source0.tex lines 547--549). The article's complete original proof of it is delivered with it (source0.tex lines 551--604, one \texttt{proof} environment of 54 lines). No mathematics is written, repaired or re-derived: the statement and the proof are the article's own lines, in the article's own notation, and no retained line is substituted or re-flowed.  Retained uncounted as the source-local context the proof needs: lines 509--511.  Nothing was omitted from any retained span, so the package carries no omission record.  No retained line contains a citation, so the wrapper carries no bibliography and no bibliography entry.  No retained line contains a figure environment, an image-inclusion command or drawing code, so no image is delivered and the package declares no asset dependency.  Editorial formatting only, in addition: an AMS \texttt{amsart} preamble that restates the article's own declarations and macros used by the retained lines, namely the environments \texttt{lemma}, \texttt{theorem} on separate counters, together with the standard packages those lines need.  The retained environments print in this document as Lemma 1, Theorem 1 in order of appearance, whereas the article prints the same items with the numbering of its own full document.  The complete native source of the article as distributed in the arXiv e-print for this exact version is bundled unchanged as \texttt{sources/196339/source0.tex}.
\begin{lemma}\label{lem:lin_basis_basis}
Let $M=(E,\rho)$ be a $q$-matroid. Any linear basis of $E$ contains a linear basis of a basis of $M$.
\end{lemma}

\begin{theorem}\label{thm:free_product_transversal}
The $q$-matroid $M_1\sq M_2$ is a transversal $q$-matroid with presentation $\mathcal{X}_1\sq\mathcal{X}_2$.
\end{theorem}

\begin{proof}
Let $M$ be the transversal $q$-matroid with presentation $\mathcal{X}_1\sq\mathcal{X}_2$. We show that $\mathcal{I}(M)=\mathcal{I}(M_1\sq M_2)$. Let $X\le E_1\oplus E_2$.

It is straightforward to verify that $M|(E_1\oplus0)\equiv M_1$ under the isomorphism $\pi_1|_{(E_1\oplus0)}$. Therefore $\pi_1(X\cap(E_1\oplus0))\in\mathcal{I}(M_1)$ is a necessary condition for $X\in\mathcal{I}(M)$, which coincides with the corresponding necessary condition for $X\in\mathcal{I}(M_1\sq M_2)$.

Given this observation, we show that $X$ is a transversal of $\mathcal{X}_1\sq\mathcal{X}_2$ if and only if
\[
k-\dim(X\cap(E_1\oplus0))
\ge
\dim(\pi_2(X))-\rho_2(\pi_2(X)).
\]

First observe that for any $v\in E_1\oplus E_2$ and $i\in[\ell]$ we have
\[
v\notin E_1\oplus B_i
\quad\text{if and only if}\quad
\pi_2(v)\notin B_i.
\]

Let $\beta$ be a linear basis of $X$. If $\beta$ admits an avoidance transversal of $\mathcal{X}_1\sq\mathcal{X}_2$, then any linear basis $\tilde{\beta}$ of $X$ satisfying
\[
\tilde{\beta}\cap(E_1\oplus0)\subseteq\beta\cap(E_1\oplus0)
\quad\text{and}\quad
\beta\setminus(E_1\oplus0)\subseteq\tilde{\beta}\setminus(E_1\oplus0)
\]
also admits such a transversal. Hence we may assume
\[
|\beta\cap(E_1\oplus0)|=\dim(X\cap(E_1\oplus0)),
\]
which implies
\[
|\beta\setminus(E_1\oplus0)|=|\pi_2(\beta)|=\dim(\pi_2(X)).
\]

Let $\beta_1=\beta\cap(E_1\oplus0)$. Since $\pi_1(X\cap(E_1\oplus0))\in\mathcal{I}(M_1)$, we may fix an avoidance transversal for $\beta_1$ using $|\beta_1|=\dim(X\cap(E_1\oplus0))$ members of $(A_1\oplus0,\dots,A_k\oplus0)$.
Since $\langle\pi_2(\beta)\rangle=\pi_2(X)$, Lemma~\ref{lem:lin_basis_basis} implies that there exists a linearly independent set $\beta'\subseteq\pi_2(\beta)$ such that $\langle\beta'\rangle$ is a basis of $\pi_2(X)$ in $M_2$. In particular $|\beta'|=\rho_2(\pi_2(X))$.
Choose $\beta_2\subseteq\beta\setminus(E_1\oplus0)$ such that $\pi_2(\beta_2)=\beta'$ and $|\beta_2|=|\beta'|$. The avoidance transversal of $\beta'$ in $\mathcal{X}_2$ induces an avoidance transversal of $\beta_2$ in $(E_1\oplus B_1,\dots,E_1\oplus B_\ell)$.

The remaining vectors of $\beta$ are
\[
\beta\setminus(\beta_1\cup\beta_2),
\]
whose number equals
\[
\dim(\pi_2(X))-\rho_2(\pi_2(X)).
\]
Thus there are enough remaining members of $(A_1\oplus0,\dots,A_k\oplus0)$ to complete an avoidance transversal of $\beta$ if and only if
\[
\dim(\pi_2(X))-\rho_2(\pi_2(X))
\le
k-\dim(X\cap(E_1\oplus0)).
\]
This shows that $X\in\mathcal{I}(M)$ if and only if $X\in\mathcal{I}(M_1\sq M_2)$, and hence $\mathcal{I}(M)=\mathcal{I}(M_1\sq M_2)$.
\end{proof}

\end{document}
